\chapter*{How to Use This Book}
\addcontentsline{toc}{chapter}{How to Use This Book}
The book has five parts. Choose a path based on your goal:
\begin{description}
\item[Part I --- Exam-Ready Reference] a compact study guide
(\cref{ch:studyguide}). If the exam is close and you only need to refresh the
\emph{circuit-theory} formulas, this part (with the practice in Part~V) is the
fastest route. The scope note in the preface explains that this book covers one
portion of the exam, not all of it.
\item[Part II --- Mathematical and Control-Theory Foundations] the generic
toolkit, with no circuits assumed: complex numbers and phasors; modeling first- and
second-order linear systems (ODEs, state space, Laplace, block diagrams); frequency
response, Bode plots, poles, zeros, and the s-plane; feedback; and how all of it
extends to higher order. Read this if you are rusty on control theory.
\item[Part III --- Circuit Models] the simple circuits, built one at a time as
instances of the Part~II toolkit: components and modeling, series and parallel
networks, AC steady state, and then RC, RL, series RLC, and parallel RLC. Here the
Part~II models become circuit behavior.
\item[Part IV --- Applying the Circuit Models] what those simple circuits explain:
filters and matching networks, transmission lines and antennas, active circuits and
oscillators, sampling and digital signal processing, receiver noise and dynamic
range, measurement, antenna patterns and practical design (\cref{ch:antennapatterns}),
signals and receiver architecture (\cref{ch:signals}), and device bias and regulation
(\cref{ch:devices}). These chapters reuse the earlier models rather than derive each
topic from scratch. \Cref{ch:dsp}, for example, starts from mixer multiplication;
\cref{ch:noise} starts from one noise formula and one power series.
\item[Part V --- Practice and Study] an exam map (\cref{ch:exammap}), reasoning
through real question-pool items (\cref{ch:practice}), and cross-chapter problems
that each require several chapters at once (\cref{ch:crossproblems}).
\end{description}

Parts~III and~IV express the book's central claim: control-theory tools describe
simple ideal circuits, and those circuits make filters, antennas, transmission
lines, and the rest of the exam's circuit content intelligible rather than arbitrary.

Each of the four circuit chapters in Part~III (RC, RL, series RLC, parallel RLC)
follows the same recipe:
\begin{enumerate}
\item draw the schematic, name the energy-storage state variables, and say what the
output is;
\item apply KCL/KVL and the element laws to get the differential equation;
\item write the state-space model and identify the poles;
\item match it to the generic first- or second-order system of Part~II;
\item interpret the frequency response (Bode) and pole motion (root locus);
\item close with the energy balance and a worked example.
\end{enumerate}

Throughout, colored sidebars separate the layers. \emph{Exam Relevance} and
\emph{Common Mistake} carry exam-level facts; \emph{Controls Connection},
\emph{Physical Insight}, and \emph{Energy Insight} carry the deeper interpretation;
\emph{Units Check} and \emph{Advanced Note} are asides; and \emph{Worked Example}
holds a calculation done in full. For a fast review, read the \emph{Exam Relevance}
and \emph{Common Mistake} boxes first. The appendices add a formula index with a
column saying where each result is derived (\cref{app:formulas}), a units and prefix
reference (\cref{app:units}), and a glossary (\cref{app:glossary}).

\paragraph{About the figures.}
\emph{Schematics} are drawn inline in the \LaTeX{} source; they are diagrams, not
measurements. Every other figure, including a Bode plot, step response, root locus,
pole diagram, Smith chart, or attenuation curve, is \emph{computed} by a same-named
Python script in \texttt{figures/src/}. For example,
\texttt{bode\_second\_order.py} generates \texttt{bode\_second\_order.pdf}.
Nothing is traced, sketched from memory, or copied from another source. A reader can
therefore regenerate or re-plot any curve by editing one short file. Captions name a
model or fitted values only when that choice matters.
